\section{Introduction}
\label{sec:intro}

Tree ensembles and regularized linear models have long been the practical default for tabular data, and extensive
comparisons found gradient-boosted trees competitive with or better than deep networks on typical medium-sized tables
\citep{shwartzziv2022tabular,grinsztajn2022tree,borisov2024survey}. Tabular foundation models take a different
approach: a transformer is pretrained once, mostly on synthetic tables, and then predicts the labels of new rows by
conditioning on the labelled rows of a new dataset, without any training on it
\citep{hollmann2023tabpfn,hollmann2025tabpfn,qu2025tabicl}. NVIDIA Kumo Tabular \citep{nvidia_kumo_card,nvidia_kumo_blog}
is a recently released model of this kind. Its developers report leading results on public benchmark suites, but we know
of no independent evaluation.

This paper reports an empirical evaluation of Kumo Tabular on small, fully numerical tables that is independent of
the model's developers. It does not propose a new method. We ask:
\begin{enumerate}
\item[\textbf{RQ1}] How does Kumo Tabular, used with default settings, compare with seven untuned classical models in
predictive quality across an externally defined set of small datasets?
\item[\textbf{RQ2}] Is any difference consistent across datasets, robust to dependence between related datasets, and
present in ranking quality rather than only in probability calibration?
\item[\textbf{RQ3}] Where does Kumo Tabular fail to outperform the classical models?
\item[\textbf{RQ4}] What does it cost in runtime on modest hardware?
\end{enumerate}

Our contributions are:
\begin{itemize}
\item A benchmark of Kumo Tabular against logistic or linear regression, random forests, extra trees, histogram
gradient boosting, XGBoost, LightGBM and CatBoost on 21 OpenML datasets selected by fixed criteria, using the official
OpenML splits (Section~\ref{sec:design}).
\item A planned analysis with datasets as the unit of inference, a single primary metric per task, Holm-corrected
pairwise tests, effect sizes and two sensitivity analyses, with the statistics recomputed by a second, separately written
implementation (Section~\ref{sec:results}).
\item Two detailed case studies, including a training-size study and an evaluation of predictive intervals, which
illustrate how dataset-level conclusions can differ from the aggregate picture (Section~\ref{sec:case}).
\item Code, raw per-fold results, logs and analysis scripts as anonymized supplementary material.
\end{itemize}
